% Add "answers" option to print the answers
%\documentclass[addpoints,answers]{exam}
\documentclass[addpoints]{exam}

\usepackage[utf8]{inputenc}
\usepackage{calc}
\usepackage{listings}
\usepackage{xcolor}
\usepackage{xspace}

\definecolor{codegreen}{rgb}{0,0.6,0}
\definecolor{codegray}{rgb}{0.5,0.5,0.5}
\definecolor{codepurple}{rgb}{0.58,0,0.82}

\lstdefinestyle{Cstyle}{
    commentstyle=\color{codegreen},
    keywordstyle=\color{magenta},
    numberstyle=\tiny\color{codegray},
    stringstyle=\color{codepurple},
    basicstyle=\ttfamily\small,
    breakatwhitespace=true,         
    breaklines=true,                 
    captionpos=b,                    
    keepspaces=true,                 
    numbers=left,                    
    numbersep=5pt,                  
    showspaces=false,                
    showstringspaces=false,
    showtabs=false,                  
    tabsize=4
}

\lstset{style=Cstyle}

% \makebox[\textwidth]{Name:\enspace\hrulefill}
%  
% \vspace{5mm}
%\makebox[\textwidth]{Instructor’s name:\enspace\hrulefill}

\newcommand{\truefalse}
{		\begin{checkboxes}
			 \choice True
			 \choice False. Explanation: 
		\end{checkboxes}
		\vspace{10pt}
}

\newcommand{\answerspace}[1]
{
	\ifprintanswers
	\else
		\setlength{\mylength}{10pt}%
		%\noindent\rule{\mylength}{20pt}
		%\vspace{\fpeval{1.2*(#1)}}
		\vspace{#1\mylength}
	\fi
}
\newlength{\mylength}

\newcommand\pf[1]{\textcolor{red}{PF: #1}}
\newcommand\am[1]{\textcolor{blue}{AM: #1}}

\newcommand{\answer}[1]{\begin{solution}{#1}\end{solution}}
\newcommand{\ra}{$\rightarrow$\xspace}



%\usepackage{listings}
\usepackage{color}
\usepackage{lastpage}
\usepackage{verbatim}
\usepackage{fancyvrb}
\usepackage{graphicx}
\usepackage{float}
\usepackage{algorithm}
\usepackage{algpseudocode}
\usepackage{minted}

%\pagestyle{fancy}
%\fancyfoot[C]{{\thepage} of \pageref{LastPage}}
\RecustomVerbatimCommand{\VerbatimInput}{VerbatimInput}{}

\renewcommand{\solutiontitle}{\noindent\textbf{Solution sketch: }}
\renewcommand{\pf}[1]{}

\newcommand{\boxedanswer}[2]{\ifprintanswers\answer{#2}\else\makeemptybox{#1}\fi}

\definecolor{SolutionColor}{rgb}{0.8,0.9,1}


\cfoot{Page \thepage\ of \pageref{LastPage}}

\rhead{Initials: \rule{20mm}{0.4pt}}



%%%%%%%%%%%%%%%%%%%%%%%%%%%%%%%%%%%%%%%%
\begin{document}
\fvset{numbers=left,numbersep=3pt}
\thispagestyle{empty}

\begin{center}
  {\LARGE CS 408 Fall 2026 -- Homework 3\\}

  {\Large Deadline: Friday, October 2nd \textbf{at 6p} \\}

  \vspace{20mm}
  \fbox{
    \fbox{\parbox{5.5in}{\centering
        \begin{itemize}
        \item The homework is open notes, open book.
        \item Please upload an annotated version of this pdf, ensuring
          that each answer fits in the box following each
          question.
        \item All solutions must be \textbf{handwritten}; typed
          work will not receive any credit. You may annotate the pdf
          directly, no need to print out and rescan.
        \item No AI tools allowed.
        \item No collaboration with your peers of any kind allowed on
          this homework.
        \item Homeworks are due at \textbf{6PM}. Anything submitted
          after \textbf{6PM} will not be accepted.
        \item Provide complete and clear but concise answers. Clearly
          explain your answers where required.
	\end{itemize}
      }}}
\end{center}

\vspace{13mm}
\begin{center}
	\makebox[0.6\textwidth]{
		\combinedgradetable[h][questions]
	}
\end{center}
%\begin{center}\combinedgradetable[h][questions]\end{center}
\vspace{28mm}

\begin{center}
\makebox[\textwidth]{\textbf{Name (please print)}: \rule{100mm}{0.8pt} }
\end{center}

\newpage
\begin{questions}

\question Consider the following function, which accepts a booking
  during business hours, i.e., between 9:00a and 5:00p, or from an
  administrator at any time.

\begin{minted}{python}
def book_seat(hour: int, is_admin: bool) -> str:
    if (hour >= 9 and hour < 17) or is_admin:
        return "booked"
    return "declined"
\end{minted}

Write $p$ for the predicate on line 2, and let
$c_1 = \texttt{hour >= 9}$, $c_2 = \texttt{hour < 17}$, and
$c_3 = \texttt{is\_admin}$.

\begin{parts}

\part[2] Write out the full truth table for $p$. Include one column for
  each clause and one for the value of $p$, and number the rows so you
  can refer to them below. Your table should have the form:
  \begin{center}
  \begin{tabular}{|c|c|c|c|c|}
    \hline
    & $c_1$ & $c_2$ & $c_3$ & $p$ \\
    \hline
    1 & $\top$ & $\top$ & $\top$ & $\top$ \\
    \hline
    2 & $\ldots$ & $\ldots$ & $\ldots$ & $\ldots$ \\
  \end{tabular}
\end{center}

  \boxedanswer{4in}{ %
    \begin{tabular}{c|ccc|c}
      & $c_1$ & $c_2$ & $c_3$ & $p$ \\
      \hline
      1 & $\top$ & $\top$ & $\top$ & $\top$ \\
      2 & $\top$ & $\top$ & $\bot$ & $\top$ \\
      3 & $\top$ & $\bot$ & $\top$ & $\top$ \\
      4 & $\top$ & $\bot$ & $\bot$ & $\bot$ \\
      5 & $\bot$ & $\top$ & $\top$ & $\top$ \\
      6 & $\bot$ & $\top$ & $\bot$ & $\bot$ \\
      7 & $\bot$ & $\bot$ & $\top$ & $\top$ \\
      8 & $\bot$ & $\bot$ & $\bot$ & $\bot$ \\
    \end{tabular}
  }

\part[1] Some of the rows in your truth table do not correspond to
  \textit{any} execution of \mintinline{python}{book_seat}. What are
  they, and why? %
  \boxedanswer{1.65in}{ %
    Rows 7 and 8, in which $c_1 = c_2 = \bot$. Both clauses read the
    same variable: $c_1 = \bot$ means \texttt{hour < 9} and
    $c_2 = \bot$ means \texttt{hour >= 17}, and no integer satisfies
    both. The remaining six rows are all reachable.  }

  \newpage
\part[3] For each of the clauses in $p$, give concrete inputs to
  \mintinline{python}{book_seat} in which that clause determines $p$.
  For each one, note which row the inputs correspond to. Your answer
  should have the form:
  \begin{itemize}
  \item[$c_i$:] \mintinline{python}{book_seat(}$n$, $b$\mintinline{python}{)}, row $j$.
  \end{itemize}
  \boxedanswer{2.5in}{ %
    Any determining row is acceptable. One convenient choice uses
    \mintinline{python}{is_admin=False} throughout and varies only
    \texttt{hour}:

    \begin{description}
    \item[$c_1$:] \mintinline{python}{book_seat(12, False)} --- row 2,
      $p = \top$. Changing \texttt{hour} to 6 flips $c_1$ and gives
      $p = \bot$.
    \item[$c_2$:] \mintinline{python}{book_seat(20, False)} --- row 4,
      $p = \bot$. Changing \texttt{hour} to 12 flips $c_2$ and gives
      $p = \top$.
    \item[$c_3$:] \mintinline{python}{book_seat(6, False)} --- row 6,
      $p = \bot$. Changing \texttt{is\_admin} to \texttt{True} gives
      $p = \top$.
    \end{description}

    In full: $c_1$ determines $p$ in rows 2 and 6, $c_2$ in rows 2 and
    4, and $c_3$ in rows 3, 4, 5 and 6 (and, were they feasible, 7 and
    8).

    Note that all three tests above hold \texttt{is\_admin} fixed and
    change only \texttt{hour}: 12, 20 and 6. Because $c_1$ and $c_2$
    read the same variable, the value of \texttt{hour} alone decides
    which clauses are in a position to determine $p$.
  }

\part[2] Looking at the truth table you provided above, you should
  observe that $c_1$ and $c_2$ can only determine $p$ for one
  particular value of \mintinline{python}{is_admin}. In two to three
  sentences, explain why this is consistent with the natural language
  description of \mintinline{python}{book_seat}'s behavior. %

  \boxedanswer{3in}{ %
    $c_1$ determines $p$ only when $c_2 = \top$ and $c_3 = \bot$, and
    $c_2$ only when $c_1 = \top$ and $c_3 = \bot$. As soon as
    $c_3 = \top$ the disjunction is satisfied, so $p = \top$ no matter
    what \texttt{hour} is --- the administrator override masks the
    business-hours check entirely.

    The same effect explains row 1: with
    \mintinline{python}{book_seat(12, True)} the booking succeeds, but
    \emph{no} clause determines $p$, because flipping any one clause
    on its own leaves it true. A test can execute a predicate and
    still tell you nothing about which clause was responsible for the
    outcome.
  }

\end{parts}

\newpage
\question Answer the following questions for the
  \mintinline{python}{checkIt} function shown below.

\begin{minted}{python}
def checkIt(a: bool, b: bool, c: int, d: int) -> None:
    if (a and b) or (c < d):
        print("P is true")
    else:
        if d < c:
            print("P isn't true and c is not equal to d")
        else:
            print("P isn't true and c is equal to d")
\end{minted}

\begin{parts}
\part[3] List all predicates in \mintinline{python}{checkIt}, followed
  by a list of the individual clauses in each predicate.
  Use the format ``predicate1: clause1, clause2, ...''
  \boxedanswer{1.5in}{ %
    $(a \wedge b) \vee (c < d): a, b, c < d$ \\
    $ d < c: d < c$
}

% One point for each in TR
\part[8] Derive a test set satisfying CACC but not RACC for the first
  predicate in \mintinline{python}{checkIt}, and explain why it fails
  RACC. %
  \boxedanswer{4.5in}{ %
	% The truth table of the predicate in (a) is: \\
	% \begin{tabular}{l|l|l|l|l}
	% 	& a & b & c & P \\ \hline
	% 1 & F & F & F & F \\ \hline
	% 2 & F & F & T & F \\ \hline
	% 3 & F & T & F & F \\ \hline
	% 4 & F & T & T & F \\ \hline
	% 5 & T & F & F & F \\ \hline
	% 6 & T & F & T & T \\ \hline
	% 7 & T & T & F & T \\ \hline
	% 8 & T & T & T & T
	% \end{tabular}
	The following forms of ACC: GACC, CACC, and RACC were accepted as correct answers.

	Make $a$ major then we need $c < d=\bot$ and $b=\top$ for $a$ to determine \(p = (a \wedge b) \vee (c < d)\).
	\[
	TR_a = \{ (a = \top, b = \top, c < d = \bot),\; (a = \bot, b = \top, c < d = \bot)\}
	\]

	Make $b$ major then we need $c=\bot$ and $a=\top$ to determine
	\(p = (a \wedge b) \vee (c < d)\).

	\[
	TR_b = \{ (a = \top, b = \bot, c < d = \bot),\; (a = \top, b = \top, c < d = \bot)\}
	\]

	Make $c$ major then we need $(a \wedge b)=\bot$ to determine
	\(p = (a \wedge b) \vee (c < d)\).

	\[
	TR_{c < d} = \{ (a = \bot, b = \bot, c < d = \bot),\; (a = \bot, b =
        \bot, c < d = \top)\}
	\]

	This gives us

	\[
	TR = TR_a \cup TR_b \cup TR_{c < d} \cup TR_{d < c}
	\]

	\[
	TR = \left\{
	\begin{aligned}
	& (a = \top, b = \top, c < d = \bot),\\
	& (a = \bot, b = \top, c < d = \bot),\\
	& (a = \top, b = \bot, c < d = \bot),\\
	& (a = \bot, b = \bot, c < d = \bot),\\
	& (a = \bot, b = \bot, c < d = \top) \\
        & (d < c = \top) \\
        & (d < c = \bot) \\
      \end{aligned}
    \right\}
  \]
  Final test set:
  \begin{description}
  \item[R1] [a $\mapsto$ true, b $\mapsto$ true, c $\mapsto$ 5, d $\mapsto$ 3]
  \item[R2] [a $\mapsto$ false, b $\mapsto$ true, c $\mapsto$ 5, d $\mapsto$ 3]
  \item[R3] [a $\mapsto$ true, b $\mapsto$ false, c $\mapsto$ 5, d $\mapsto$ 3]
  \item[R4] [a $\mapsto$ false, b $\mapsto$ false, c $\mapsto$ 2, d $\mapsto$ 5]
  \item[R5] [a $\mapsto$ false, b $\mapsto$ false, c $\mapsto$ 5, d $\mapsto$ 2]
  \end{description}
}

\end{parts}

\newpage

\question The function \mintinline{python}{average_in_range()} below
  is intended to return the average (rounded down) of the elements of
  \mintinline{python}{arr} that are strictly positive and strictly
  less than \mintinline{python}{cap}, and to return 0 when there are
  no such elements.

\begin{minted}{python}
def average_in_range(arr, cap):
    total = 0
    count = 0
    i = 0
    while i < len(arr):
        if arr[i] > 0 and arr[i] < cap:
            total = total + arr[i]
            count = count + 1
        i = i + 1
    return total // count
\end{minted}

  Throughout this question, consider the test set
  \begin{description}
    \item[T$_1$:] \mintinline{python}{average_in_range([5, -1, 20], 10)}
    \item[T$_2$:] \mintinline{python}{average_in_range([4, 6], 10)}
    \end{description}
    Note that both of these tests return the expected value, and thus
    pass.

\begin{parts}

\part[3] Draw the control flow graph of
  \mintinline{python}{average_in_range()}, grouping statements into
  basic blocks and labelling each node with the line numbers it
  covers.
  \boxedanswer{5in}{ %
    Six nodes: \texttt{BB2-4} (the three initializations), \texttt{C5}
    (the loop condition), \texttt{C6} (the \texttt{if} condition),
    \texttt{BB7-8} (the two updates), \texttt{BB9}
    (\texttt{i = i + 1}), and \texttt{BB10} (the return).

    Seven edges: \texttt{(BB2-4, C5)}, \texttt{(C5, C6)},
    \texttt{(C5, BB10)}, \texttt{(C6, BB7-8)}, \texttt{(C6, BB9)},
    \texttt{(BB7-8, BB9)}, \texttt{(BB9, C5)}.
  }
\part[1] How many test requirements does node coverage impose?
  \boxedanswer{1in}{ %
    Node coverage imposes 6 requirements.
  }

\part[1] How many does edge coverage impose?
  \boxedanswer{1in}{ %
    Edge coverage imposes 7 requirements.
  }

\part[2] Give the test paths toured by each of the two tests
  above. Taken together, do they achieve 100\% edge coverage? %
  \boxedanswer{2in}{ %
    $T_1$: \texttt{[BB2-4, C5, C6, BB7-8, BB9, C5, C6, BB9, C5, C6,
      BB9, C5, BB10]}

    $T_2$: \texttt{[BB2-4, C5, C6, BB7-8, BB9, C5, C6, BB7-8, BB9, C5,
      BB10]}

    $T_1$ alone tours all seven edges: element 5 takes \texttt{(C6,
      BB7-8)}, elements $-1$ and 20 both take \texttt{(C6, BB9)}, and
    the final loop test takes \texttt{(C5, BB10)}. So this pair of
    tests achieves 100\% edge coverage (and therefore 100\% node
    coverage).  }

\part[4] Show that this set of tests satisfies Correlated Active
  Clause Coverage for the predicate on line 6 by identifying which
  required truth assignment each test supplies, noting the major and
  minor clauses in each requirement. %
  \boxedanswer{2.5in}{ %
    Write $c_1$ for \mintinline{python}{arr[i] > 0} and $c_2$ for
    \mintinline{python}{arr[i] < cap}, so $p = c_1 \wedge c_2$.

    \textbf{$c_1$ major:} $c_1$ determines $p$ when $c_2 = \top$. Need
    $(c_1, c_2) = (\top, \top)$ with $p = \top$ and $(\bot, \top)$
    with $p = \bot$. In $t_1$: element 5 gives $(\top, \top)$; element
    $-1$ gives $(\bot, \top)$, since $-1 < 10$.

    \textbf{$c_2$ major:} $c_2$ determines $p$ when $c_1 = \top$. Need
    $(\top, \top)$ and $(\top, \bot)$. In $t_1$: element 5 gives
    $(\top, \top)$; element 20 gives $(\top, \bot)$.

    All four evaluations occur within $t_1$, so $T$ satisfies CACC.
  }
\newpage
\part[5] List all du-paths with respect to \mintinline{python}{count}.
  For each one, state whether this test set tours it.
  \boxedanswer{4in}{ %
    \mintinline{python}{count} is defined at \texttt{BB2-4} and at
    \texttt{BB7-8}, and used at \texttt{BB7-8} and \texttt{BB10}. Four
    du-paths:

    \begin{tabular}{llc}
      \textbf{du-pair} & \textbf{du-path} & \textbf{toured by $T$?} \\
      \hline
      (BB2-4, BB7-8) & [BB2-4, C5, C6, BB7-8] & yes ($t_1$, $t_2$) \\
      (BB2-4, BB10) & [BB2-4, C5, BB10] & \textbf{no} \\
      (BB7-8, BB7-8) & [BB7-8, BB9, C5, C6, BB7-8] & yes ($t_2$) \\
      (BB7-8, BB10) & [BB7-8, BB9, C5, BB10] & yes ($t_2$) \\
    \end{tabular}

    The second one is untoured because both tests enter the loop body
    at least once, so neither goes straight from \texttt{C5} to
    \texttt{BB10}.
  }

\part[2] Does this test set satisfy all-defs coverage with respect to
  \mintinline{python}{count}? Does it satisfy all-uses coverage?
  Justify both answers in one to two sentences.  %
  \boxedanswer{3in}{ %
    \textbf{All-defs: yes.} Each def needs only one du-path toured.
    The def at \texttt{BB2-4} is covered by \texttt{[BB2-4, C5, C6,
      BB7-8]}, and the def at \texttt{BB7-8} by \texttt{[BB7-8, BB9,
      C5, BB10]}.

    \textbf{All-uses: no.} Every (def, use) pair needs its own
    du-path, and the pair (\texttt{BB2-4}, \texttt{BB10}) is not
    toured by either test.
  }

  \newpage
\part[3] Neither test induces a failure, but
  \mintinline{python}{average_in_range()} nevertheless contains a
  fault. Give a test input that reveals it, and explain why satisfying
  all-uses coverage with respect to \mintinline{python}{count} would
  have forced such a test. %
  \boxedanswer{4in}{ %
    The function never guards against \mintinline{python}{count == 0},
    so line 10 divides by zero.
    \mintinline{python}{average_in_range([], 10)} raises
    \texttt{ZeroDivisionError} where the specification requires 0.

    Touring \texttt{[BB2-4, C5, BB10]} means leaving the loop without
    ever entering the body, which is only possible when
    \mintinline{python}{arr} is empty --- and that is exactly the case
    in which \mintinline{python}{count} is still 0 at line 10.

    (Edge-pair coverage would also force it, since
    \texttt{[BB2-4, C5, BB10]} is a path of length 2. Note also that
    \mintinline{python}{average_in_range([-1], 10)} fails for the same
    reason, but no criterion in this question requires it.)
  }

\part[3] In one or two sentences, what does this example show about the
  relationship between a test suite's coverage and its ability to
  detect faults?
  \boxedanswer{2.5in}{ %
    Satisfying a criterion guarantees only that certain paths are
    executed, not that a fault is exposed --- $T$ reaches every node,
    every edge, and every clause-determining evaluation, and still
    misses the fault. Stronger criteria shrink the gap but never close
    it, and no criterion says anything about whether the oracle
    actually checks the right thing.
  }

\end{parts}

\end{questions}

\end{document}
